\begin{abstract}
Agents are widely deployed as assistants over documents, tools, and code. However, they typically act only on explicit user requests, which surface only the problems the user has noticed, while many other important problems coexist, hidden in plain sight, within the broader user context, with their total number unknown in advance. We frame this as the task of \textit{discovering multiple hidden problems from context}, in which coexisting problems should be uncovered, grounded in supporting evidence, and paired with concrete actions. To this end, we introduce \textsc{TIDE}, a template-guided iterative framework with two complementary mechanisms. Specifically, motivated by the observation that single-pass prediction anchors on the most salient cases and yields generic claims, we propose \textit{iterative discovery}, which surfaces a small batch of candidates per round while conditioning on what has already been found, so subsequent rounds extend coverage; and \textit{thought templates}, reusable schemas distilled from previously solved cases that specify what contextual signals to attend to and how to connect them, anchoring each prediction in a recognizable problem class. We validate \textsc{TIDE} on two realistic settings, personal workspaces and software repositories, across four model backbones, showing substantial gains over single-shot and parallel multi-agent baselines on task coverage, identification, and resolution.
\end{abstract}
